\documentclass[11pt,a4paper]{article}
\usepackage{fontspec}
\setmainfont{Latin Modern Roman}
\setmonofont{Latin Modern Mono}[Scale=MatchLowercase]
\usepackage{unicode-math}
\setmathfont{Latin Modern Math}
\usepackage{newunicodechar}
\newfontfamily\grk{Liberation Serif}[Scale=MatchLowercase]
\newunicodechar{β}{{\grk β}}
\newunicodechar{−}{\ensuremath{-}}
\usepackage[a4paper,margin=1in]{geometry}
\usepackage{graphicx}
\usepackage{microtype}
\usepackage[hidelinks]{hyperref}
\setcounter{secnumdepth}{1}
\setcounter{tocdepth}{1}
\setlength{\emergencystretch}{3em}
\sloppy
\setlength{\parskip}{0.5\baselineskip}
\setlength{\parindent}{0pt}
\newcommand{\sfig}[5]{%
  \phantomsection\addcontentsline{toc}{section}{Supplementary Figure #4}%
  \null\vfill
  \begin{center}\includegraphics[width=#5\textwidth,height=0.74\textheight,keepaspectratio]{#1}\end{center}
  \medskip
  {\normalsize\bfseries\hyphenpenalty=10000\exhyphenpenalty=10000 Supplementary Figure #2\par}\smallskip
  {\small #3\par}
  \vfill
  \clearpage}
\newcommand{\sfigblk}[5]{%
  \phantomsection\addcontentsline{toc}{section}{Supplementary Figure #4}%
  \begin{center}\includegraphics[width=#5\textwidth,height=0.50\textheight,keepaspectratio]{#1}\end{center}
  \medskip
  {\normalsize\bfseries\hyphenpenalty=10000\exhyphenpenalty=10000 Supplementary Figure #2\par}\smallskip
  {\small #3\par}}
\newcommand{\sfigside}[5]{%
  \phantomsection\addcontentsline{toc}{section}{Supplementary Figure #4}%
  \noindent
  \begin{minipage}[c]{#5\textwidth}\centering\includegraphics[width=\linewidth]{#1}\end{minipage}%
  \hfill
  \begin{minipage}[c]{\dimexpr0.94\textwidth-#5\textwidth\relax}%
    {\normalsize\bfseries\hyphenpenalty=10000\exhyphenpenalty=10000 Supplementary Figure #2\par}\smallskip
    {\small #3\par}%
  \end{minipage}}
\begin{document}
{\Large\bfseries Supplementary Figures\par}\vspace{8pt}
{\bfseries EMToolbox uncovers biased co-option of embryonic EMT programs across cancers\par}\vspace{2pt}
Caleb C. Reagor, CO-AUTHORS, and Igor Adameyko\par\vspace{14pt}
{\setlength{\parskip}{0pt}\tableofcontents}
\clearpage
\sfig{figure-s1.pdf}{1: STRING network for Gast}{STRING protein-protein interaction network for Gast. Dashed outlines indicate functional enrichments, including Wnt, Notch, and TGF-β signaling.}{1: STRING network for Gast}{1}
\sfig{figure-s2.pdf}{2: STRING network for pSom}{STRING protein-protein interaction network for pSom. Dashed outlines indicate functional enrichments, including Wnt and retinoic acid signaling.}{2: STRING network for pSom}{1}
\sfig{figure-s3.pdf}{3: STRING network for Crest}{STRING protein-protein interaction network for Crest. Dashed outlines indicate functional enrichments, including Wnt and retinoic acid signaling.}{3: STRING network for Crest}{1}
\sfigblk{figure-s4.pdf}{4: Gene expression magnitude and timing correspond across e-EMT trajectories}{\textbf{(a, b)} Scatter plots depict gene expression magnitude (a) and timing (b) for the 237 genes comprising MesenCoder's input features, compared pairwise across the gastrulation, posterior somitogenesis, and neural crest migration trajectories. Labels denote genes with the largest absolute residuals from the OLS fits (red lines), and correlations (red annotations) encompass all plotted genes.}{4: Gene expression magnitude and timing}{0.95}
\par\vfill
\sfigside{figure-s5.pdf}{5: HNSCC sample integration}{UMAP embedding of malignant cells from head and neck squamous cell carcinoma (Puram \textit{et al.}, 2017), integrated across samples and colored by sample.}{5: HNSCC sample integration}{0.42}
\clearpage
\sfigblk{figure-s6.pdf}{6: Drug development programs have neglected e-EMT modules}{\textbf{(a)} Bar plots showing the proportion of each module's proteins with oncology target-indication pairs (left), their mean number of indications (middle), and their proportion entering clinical trials (right) (Minikel \textit{et al.}, 2024). Error bars denote the standard error of the mean. *, \textit{P} < 0.05; **, \textit{P} < 0.01. \textbf{(b)} Grouped bar plot showing the proportion of each module's proteins assigned to Pharos target development levels (Kelleher \textit{et al.}, 2023). Tdark, unstudied; Tbio, biologically characterized; Tchem, bound by potent small molecules; Tclin, targeted by an approved drug.}{6: Drug development programs}{1}
\par\vfill
\sfigblk{figure-s7.pdf}{7: Nominated targets' normal expression and pan-cancer coexpression}{\textbf{(a)} Violin plot showing the expression of the nominated targets (light blue) and \textit{GPC3} (light red) across normal cell types from postembryonic tissues (CZI Single-Cell Biology Program \textit{et al.}, 2023). Nominated targets exhibit maximum normal expression below that of \textit{GPC3}, a clinically tolerated vaccine target. Dashed lines indicate the median and upper/lower quartiles. \textbf{(b)} Dotplot showing the relative expression of nominated targets across 3CA diseases (Tyler \textit{et al.}, 2025). Vertical lines denote target coexpression.}{7: Nominated targets' expression}{1}
\par\vspace{\stretch{2}}\clearpage
\end{document}
